% !TeX program = pdflatex
\documentclass[12pt]{article}

% Idioma y codificación
\usepackage[spanish]{babel}
\usepackage[utf8]{inputenc}
\usepackage[T1]{fontenc}
\usepackage{lmodern}

% Página y espaciado
\usepackage{geometry}
\geometry{letterpaper, margin=1in}
\usepackage{setspace}
\onehalfspacing

% Tipografía y microajustes
\usepackage{microtype}
\usepackage{csquotes}

% Tablas
\usepackage{booktabs}
\usepackage{longtable}
\usepackage{array}

% Listas
\usepackage{enumitem}

% Código
\usepackage{listingsutf8}
\usepackage{listings}
\lstset{inputencoding=utf8}
\lstset{literate=%
  {á}{{\'a}}1 {é}{{\'e}}1 {í}{{\'\i}}1 {ó}{{\'o}}1 {ú}{{\'u}}1
  {Á}{{\'A}}1 {É}{{\'E}}1 {Í}{{\'I}}1 {Ó}{{\'O}}1 {Ú}{{\'U}}1
  {ñ}{{\~n}}1 {Ñ}{{\~N}}1 {ü}{{\"u}}1 {Ü}{{\"U}}1
}
\usepackage{xcolor}
\lstdefinestyle{python}{
  language=Python,
  inputencoding=utf8,
  basicstyle=\ttfamily\small,
  numbers=left,
  numberstyle=\tiny,
  stepnumber=1,
  numbersep=8pt,
  breaklines=true,
  showstringspaces=false,
  tabsize=2,
  frame=single
}

% Hipervínculos
\usepackage{hyperref}
\hypersetup{
  colorlinks=true,
  linkcolor=blue,
  urlcolor=blue,
  citecolor=blue
}

% Formato académico y paginación por sección
\usepackage{titlesec}
\usepackage{etoolbox}
\titlespacing*{\section}{0pt}{3ex plus 1ex minus .2ex}{2ex}
\titlespacing*{\subsection}{0pt}{2.5ex plus 0.8ex minus .2ex}{1.2ex}
\titlespacing*{\subsubsection}{0pt}{2ex plus 0.6ex minus .2ex}{1ex}
\preto\section{\clearpage}

% Encabezados/pies
\usepackage{fancyhdr}
\pagestyle{fancy}
\fancyhf{}
\lhead{Documento de apoyo (Sesión 6, Unidad 2)}
\rhead{Análisis y ecosistemas Big Data con Spark}
\cfoot{\thepage}
\renewcommand{\headrulewidth}{0.4pt}
\renewcommand{\footrulewidth}{0pt}

% Bibliografía
\usepackage[
  backend=biber,
  style=ieee,
  sorting=nyt,
  maxbibnames=99
]{biblatex}
\addbibresource{referencias_s6.bib}

% Metadatos
\title{\textbf{Documento de apoyo (Sesión 6, Unidad 2)}\\
\large Análisis y ecosistemas Big Data con Apache Spark\\
\normalsize Fundamentos, arquitectura, exploración y carga desde data lakes}
\author{Javier Andrés Valdez González / Auxiliar \\
\small Seminario de Sistemas 2 --- Sección P \\
\small Universidad de San Carlos de Guatemala}
\date{Marzo de 2026}

\begin{document}
\maketitle

\begin{abstract}
El análisis Big Data requiere arquitecturas distribuidas capaces de procesar grandes volúmenes de información con velocidad y flexibilidad. En esta sesión se introducen los fundamentos conceptuales de Big Data a través del marco de las 4V, y se desarrolla el papel de Apache Spark como motor de análisis distribuido para cargas batch e interactivas. El documento presenta la arquitectura de Spark, sus componentes principales, el flujo básico de trabajo para análisis de datos y prácticas de exploración con DataFrames. Finalmente, se aborda la carga de información desde data lakes, destacando formatos columnares y consideraciones de particionamiento para mejorar desempeño analítico.
\end{abstract}

\textbf{Palabras clave:} Big Data, 4V, Apache Spark, DataFrames, data lake, particionamiento, Parquet.

\newpage
\tableofcontents

\section{Introducción}
Las organizaciones modernas generan datos desde aplicaciones transaccionales, sensores, plataformas web, redes sociales y sistemas empresariales. Este crecimiento no solo incrementa el volumen de datos, sino también la diversidad de formatos y la velocidad con la que la información cambia. En ese contexto surge Big Data como un enfoque para almacenar, procesar y analizar datos que superan las capacidades de herramientas tradicionales \cite{ibm_big_data}.

Apache Spark se ha consolidado como uno de los motores más utilizados para análisis distribuido por su modelo de ejecución en memoria, su API unificada y su integración con ecosistemas de almacenamiento masivo. Spark permite construir pipelines de transformación, exploración y análisis sobre datos estructurados y semiestructurados en entornos locales, clústeres on-premise o nube \cite{spark_docs_overview}.

El objetivo de este documento es presentar, en formato técnico y práctico, los temas de la semana 6: definición de Big Data (4V), arquitectura de Spark, fundamentos de análisis en Spark, exploración de datos y carga desde data lakes.

\section{Análisis y ecosistemas Big Data: ¿Qué es Big Data? (4V)}
Big Data describe conjuntos de datos cuya escala, complejidad y dinámica requieren arquitecturas distribuidas y técnicas especializadas para su procesamiento \cite{laney_3v}. Una forma estándar de entenderlo es mediante las 4V:

\begin{enumerate}[label=\textbf{V\arabic*.}, leftmargin=*]
\item \textbf{Volumen:} cantidad masiva de datos (GB, TB, PB) provenientes de múltiples fuentes.
\item \textbf{Velocidad:} ritmo al que los datos se generan, transmiten y deben procesarse.
\item \textbf{Variedad:} coexistencia de formatos estructurados, semiestructurados y no estructurados.
\item \textbf{Veracidad:} nivel de calidad, consistencia y confiabilidad del dato para análisis.
\end{enumerate}

Estas dimensiones obligan a diseñar soluciones que combinen almacenamiento escalable, procesamiento paralelo, control de calidad y gobierno de datos.

\section{Análisis y ecosistemas Big Data: Arquitectura de Spark}
Spark opera sobre un modelo distribuido orientado a clúster, donde una aplicación se divide en tareas paralelas ejecutadas por nodos trabajadores. Sus componentes principales son:

\begin{itemize}[leftmargin=*]
\item \textbf{Driver:} proceso que contiene la lógica de la aplicación, planifica tareas y coordina la ejecución.
\item \textbf{Cluster Manager:} asigna recursos entre aplicaciones (por ejemplo, Standalone, YARN o Kubernetes).
\item \textbf{Executors:} procesos en nodos trabajadores que ejecutan tareas y almacenan datos en memoria/disco.
\item \textbf{SparkSession:} punto de entrada para trabajar con DataFrames, SQL y lectura/escritura de datos.
\end{itemize}

A nivel lógico, Spark transforma el código en un DAG (grafo acíclico dirigido) de etapas y tareas. Este plan se optimiza antes de ejecutarse para reducir movimientos de datos y mejorar rendimiento \cite{spark_sql_guide}.

\section{Análisis y ecosistemas Big Data: Fundamentos de Spark para análisis}
En análisis de datos, Spark se usa principalmente mediante DataFrames y Spark SQL. El flujo básico de trabajo suele ser:

\begin{enumerate}[label=\textbf{F\arabic*.}, leftmargin=*]
\item \textbf{Inicializar la sesión de Spark.}
\item \textbf{Cargar datos} desde archivos o tablas.
\item \textbf{Transformar} (filtrar, agregar, limpiar, enriquecer).
\item \textbf{Analizar} con operaciones SQL o API de DataFrame.
\item \textbf{Persistir resultados} en formatos analíticos.
\end{enumerate}

\begin{lstlisting}[style=python, caption={Ejemplo base de análisis con PySpark DataFrame}]
from pyspark.sql import SparkSession
from pyspark.sql.functions import col, sum as sum_

spark = SparkSession.builder \
    .appName("Semana6-BigData") \
    .getOrCreate()

df = spark.read.option("header", True).csv("data/ventas.csv")

resumen = (
    df.groupBy("categoria")
      .agg(sum_(col("monto").cast("double")).alias("ventas_totales"))
      .orderBy(col("ventas_totales").desc())
)

resumen.show(20, truncate=False)
\end{lstlisting}

\section{Análisis y ecosistemas Big Data: Exploración de datos en Spark}
La exploración inicial permite entender estructura, calidad y patrones antes de modelar. En Spark, algunas operaciones esenciales son:

\begin{itemize}[leftmargin=*]
\item \texttt{printSchema()} para validar tipos de datos.
\item \texttt{show()} para inspección rápida de registros.
\item \texttt{describe()} para estadísticas descriptivas básicas.
\item \texttt{count()}, \texttt{distinct()} y \texttt{groupBy()} para perfiles de distribución.
\item análisis de nulos y valores atípicos con reglas de negocio.
\end{itemize}

\begin{lstlisting}[style=python, caption={Rutina básica de exploración en Spark}]
from pyspark.sql.functions import col, count, when

print("Total de filas:", df.count())

df.printSchema()
df.show(10, truncate=False)

df.describe().show()

nulos = df.select([
    count(when(col(c).isNull(), c)).alias(c)
    for c in df.columns
])

nulos.show(truncate=False)
\end{lstlisting}

Una exploración adecuada evita errores posteriores de agregación, sesgo por datos faltantes y fallos de tipado en cálculos.

\section{Análisis y ecosistemas Big Data: Carga desde data lakes}
Un data lake almacena datos en bruto y curados, generalmente en almacenamiento de objetos o sistemas distribuidos, para su consumo analítico por distintos motores \cite{aws_datalake}. Spark se integra de forma natural con esta arquitectura gracias a sus conectores y soporte de formatos abiertos.

\subsection{Formatos comunes de lectura}
\begin{itemize}[leftmargin=*]
\item \textbf{CSV/JSON:} útiles para ingestión inicial, pero menos eficientes para análisis masivo.
\item \textbf{Parquet/ORC:} formatos columnares comprimidos y optimizados para consultas analíticas.
\item \textbf{Delta Lake (cuando aplica):} agrega transaccionalidad y versionado sobre archivos de data lake \cite{delta_docs}.
\end{itemize}

\subsection{Ejemplo de carga desde data lake}
\begin{lstlisting}[style=python, caption={Lectura y escritura en formato Parquet particionado}]
from pyspark.sql.functions import year, month, to_date, col

raw = spark.read.parquet("s3a://mi-datalake/raw/ventas/")

curated = (
    raw.withColumn("fecha", to_date(col("fecha")))
       .withColumn("anio", year(col("fecha")))
       .withColumn("mes", month(col("fecha")))
       .filter(col("monto").isNotNull())
)

(curated.write
    .mode("overwrite")
    .partitionBy("anio", "mes")
    .parquet("s3a://mi-datalake/curated/ventas/"))
\end{lstlisting}

El particionamiento por fecha facilita el \textit{pruning} de particiones, reduciendo el volumen leído y acelerando consultas.

\section{Buenas prácticas mínimas para la semana 6}
\begin{longtable}{@{}p{0.30\textwidth}p{0.62\textwidth}@{}}
\toprule
\textbf{Tema} & \textbf{Práctica recomendada} \\
\midrule
\endfirsthead
\toprule
\textbf{Tema} & \textbf{Práctica recomendada} \\
\midrule
\endhead

Lectura de datos & Definir esquema explícito cuando sea posible para evitar inferencia costosa o errores de tipos. \\

Transformaciones & Encadenar operaciones de forma declarativa y validar resultados parciales con muestras controladas. \\

Exploración & Revisar nulos, cardinalidad y rangos antes de construir agregaciones o métricas de negocio. \\

Data lake & Preferir Parquet para capas analíticas y organizar rutas particionadas por columnas de alta selectividad. \\

Rendimiento & Evitar acciones innecesarias repetidas; usar cache solo cuando un DataFrame se reutiliza varias veces. \\

\bottomrule
\end{longtable}

\section{Conclusiones}
Big Data se entiende mejor como un problema de ingeniería de datos definido por volumen, velocidad, variedad y veracidad. En este contexto, Spark aporta un modelo distribuido y unificado para procesar datos a escala con APIs orientadas a análisis.

Para la práctica de esta semana, los puntos críticos son: comprender la arquitectura de Spark, dominar el flujo base con DataFrames, explorar datos con criterios de calidad y cargar/organizar información desde data lakes con formatos eficientes. Estos fundamentos son la base para construir pipelines analíticos robustos en sesiones posteriores.

\newpage
\printbibliography
\end{document}
